\section{Conditioning on one exact length}\label{R:sec:exact}
The sum of all bridge masses at lengths near $n$ does not bound the mass at length $n$. Local renewal estimates belong to classical infinite-mean renewal theory; Garsia and Lamperti and Caravenna and Doney treat regularly varying tails~\cite{GarsiaLamperti1962,CaravennaDoney2019}. Our finite inputs give comparable powers, without a regular-variation limit. We therefore prove the local lower bound directly by exponential tilting and Fourier inversion, then apply it to the joint height--length--diameter law. The probability statement holds for every exponent $0<a<1$; the restriction $a>1/2$ enters only when we bound the contribution from small piece counts in the bridge application.

\subsection{A local renewal lower bound with bounded-factor tails}
\begin{lemma}[Annular mass and concentration]\label{R:exact:concentration}
Let $X$ be a positive integer-valued random variable of span one, and let $0<a<1$. Suppose there are constants $c,C>0$ such that
\[
 \Prob(X>x)\le Cx^{-a},\qquad
 c u^a\le1-\E e^{-uX}\le Cu^a
\]
for all sufficiently large $x$ and all sufficiently small $u>0$, respectively.
Then, for some fixed $M>1$ and $c_1>0$, every sufficiently large $x$ satisfies
\[
 \Prob(x/M<X\le x)\ge c_1x^{-a}.
\]
If $S_k=X_1+\cdots+X_k$ is a sum of independent copies, then, for all $k\ge1$ and $t\ge1$,
\[
 \sup_j\Prob(S_k=j)\le C_1k^{-1/a},\qquad
 \Prob(S_k>t)\le C_1 kt^{-a}.
\]
\end{lemma}
\begin{proof}
Evaluate the Laplace deficit at $A/x$. Its contribution from $X>x$ is at most $Cx^{-a}$, while its contribution from $X\le x/M$ is at most
\[
 \frac A x\E[X;X\le x/M]\le C A M^{a-1}x^{-a}.
\]
Choose $A$ large enough that the lower deficit $cA^ax^{-a}$ dominates the first term, and then $M$ large enough to dominate the second. The remaining contribution is at most $\Prob(x/M<X\le x)$, proving the claim.

Let $\phi(\theta)=\E e^{i\theta X}$. Use two independent copies in
$1-|\phi(\theta)|^2=\E[1-\cos(\theta(X-X'))]$. Keep $X'$ in a fixed bounded set of positive probability and $X$ in the annulus $[\epsilon/(M|\theta|),\epsilon/|\theta|]$, with a fixed small $\epsilon>0$. For sufficiently small $|\theta|$, the phase is bounded above by a small constant and below in magnitude by a fixed positive constant depending only on $\epsilon,M$. The annular bound therefore gives
$|\phi(\theta)|\le e^{-c|\theta|^a}$. Away from zero, span one gives a uniform bound strictly below one. Fourier inversion yields
\[
 \sup_j\Prob(S_k=j)\le(2\pi)^{-1}\int_{-\pi}^{\pi}|\phi(\theta)|^k\,d\theta
 \le C k^{-1/a}.
\]
Finally, either some $X_i>t$ or $\sum_iX_i\ind_{X_i\le t}>t$. The union bound and Markov's inequality, using $\E[X;X\le t]\le Ct^{1-a}$, give the sum-tail estimate.
\end{proof}

\begin{theorem}[Local renewal mass]\label{R:exact:local}
Under the hypotheses of Lemma~\ref{R:exact:concentration}, for every sufficiently large integer $m$,
\[
 V_m:=\sum_{k\ge1}\Prob(S_k=m)\ge c_2m^{a-1}.
\]
The constants depend only on the law and the constants in the hypotheses. No regular-variation limit of its tail is required.
\end{theorem}
\begin{proof}
For $u>0$ tilt the law by
\[
 p_u(j)=\frac{e^{-uj}\Prob(X=j)}{Z(u)},\qquad Z(u)=\E e^{-uX}.
\]
Write $\mu_u,v_u$ for its mean and variance. Tail integration gives the upper bounds for its first three moments. The annular mass at scale $u^{-1}$ gives the lower bounds for the first two moments; for the variance, apply the identity $2v_u=\E_u(X-X')^2$, keeping one variable bounded and the other in that annulus. Thus
\[
 \mu_u\asymp u^{a-1},\quad v_u\asymp u^{a-2},\quad
 \E_u|X-\mu_u|^3\le Cu^{a-3}.
\]
Let $\phi_u$ be the tilted characteristic function. The same two-copy argument, now using the annulus of size $u^{-1}$ for $|\theta|\le u$ and of size $|\theta|^{-1}$ for $u\le|\theta|\le\theta_0$, gives uniform bounds
\begin{equation}\label{R:exact:fourier-bounds}
 |\phi_u(\theta)|\le
 \begin{cases}
 e^{-cv_u\theta^2},& |\theta|\le u,\\
 e^{-c|\theta|^a},& u\le|\theta|\le\theta_0,\\
 1-c,&\theta_0\le|\theta|\le\pi.
 \end{cases}
\end{equation}
In the first annulus the tilt is bounded below by a positive constant, and the phase is comparable to $|\theta|/u$. In the second it is again bounded below because $u/|\theta|\le1$. The final line follows from span one and the total-variation convergence $p_u\to p$.

Fix a large constant $A$ and put $u=A/m$. Consider integers $k$ with
\[
 |m-k\mu_u|\le\tfrac1{10}\sqrt{k v_u}.
\]
There are at least $c_A m^a$ such integers for all large $m$. Indeed their center is $m/\mu_u\asymp_A m^a$ and their interval width is comparable to $\sqrt{(m/\mu_u)v_u}/\mu_u\asymp_A m^a$. On this interval $ku^a\asymp A$, and $\sqrt{kv_u}\asymp_A m$.

We verify a uniform local lower bound under the tilted law. Write $\sigma=\sqrt{kv_u}$. On $|\theta|\le\omega/\sigma$, Taylor expansion of the centered characteristic function gives its Gaussian approximation with error at most $C_\omega A^{-1/2}$, since
\[
 \frac{k\E_u|X-\mu_u|^3}{(kv_u)^{3/2}}
 \le \frac C{\sqrt{ku^a}}\le\frac C{\sqrt A}.
\]
After the change of variables $t=\sigma\theta$, Fourier inversion over this interval therefore differs by at most $C_\omega A^{-1/2}/\sigma$ from
\[
 \frac1{2\pi\sigma}\int_{-\omega}^{\omega}
 e^{-t^2/2}e^{-it(m-k\mu_u)/\sigma}\,dt.
\]
The latter has real part at least $c/\sigma$, uniformly for the allowed $k$, once $\omega$ is large enough. The portion $\omega/\sigma<|\theta|\le u$ is at most $C e^{-c\omega^2}/\sigma$ by the first line of \eqref{R:exact:fourier-bounds}. The portion $u<|\theta|\le\theta_0$ is at most $Cu e^{-c'A}$, by substituting $|\theta|=u t$ into the second line; relative to $1/\sigma$ this is at most $C\sqrt A e^{-c'A}$. The final portion is exponentially small in $k$. Choose first $\omega$ and then $A$ large enough. We obtain
\[
 \Prob_u(S_k=m)\ge c_A/m.
\]
Removing the tilt multiplies this by $e^{um}Z(u)^k$. Since $um=A$ and $k(1-Z(u))=O(A)$, this factor is bounded below by a positive constant depending only on $A$. Summing over the $c_A m^a$ choices of $k$ proves the theorem.
\end{proof}

\subsection{Uniform bridges at every compatible length}
\begin{theorem}\label{R:exact:bridge}
Choose a strict bridge uniformly from all bridges with $n$ visited vertices, rooted at a fixed bottom port, with terminal height and port free. For every $\xi>0$, along every sufficiently large even integer $n$,
\[
 \Prob_n\{n^{3/4-\xi}\le A\le D\le n^{3/4+\xi}\}\longrightarrow1.
\]
Here $A$ is the Euclidean distance between the ports and $D$ the diameter of the visited trace together with its ports. Replacing either by its vertex-endpoint convention changes it by a bounded amount.
\end{theorem}
\begin{proof}
Port bridge lengths are even: their first and last triangles have opposite types. The irreducible law admits two- and four-vertex zigzags in one row, so $X=L/2$ has atoms at $1$ and $2$ and has span one. Proposition~\ref{R:renew:increments} verifies the hypotheses of Theorem~\ref{R:exact:local} with $a=\alpha=9/16$. Thus, writing $m=n/2$, the total critical bridge mass at that length is
\[
 V_m=\sum_{k\ge1}\Prob_p\{\textstyle\sum_{i\le k}L_i=2m\}
 \ge c m^{\alpha-1}.
\]
All bridges of this length have the same critical weight, so normalization by $V_m$ gives the uniform law.

We first restrict the number $k$ of pieces. For $k\ge2$ split the list into two comparable halves. If the total $X$-sum is $m$, one half is at least $m/2$. Lemma~\ref{R:exact:concentration} bounds that probability by $Ckm^{-\alpha}$; the other independent half pays at most $Ck^{-1/\alpha}$ for the required residual length. Therefore, for any fixed $\delta>0$,
\[
 \sum_{2\le k\le m^{\alpha-\delta}}\Prob(S_k=m)
 \le C m^{-\alpha}\sum_{k\le m^{\alpha-\delta}}k^{1-1/\alpha}
 \le C m^{\alpha-1-\delta(2-1/\alpha)}.
\]
The exponent $2-1/\alpha=2/9$ is positive. The $k=1$ term is at most $Cm^{-\alpha}=o(m^{\alpha-1})$. At the other end, exponential Markov at discount $1/m$ gives
\[
 \Prob(S_k=m)\le e\exp(-c k m^{-\alpha}),
\]
so the sum over $k>m^{\alpha+\delta}$ is negligible. Here $k\le m$ because $X\ge1$.

For the remaining $k$, total height below $t=m^{3/4-\xi}$ costs at most
\[
 \Prob\{\textstyle\sum H_i\le t\}
 \le e\exp(-c k t^{-\beta})
 \le e\exp(-c m^{\beta\xi-\delta}),
\]
which is negligible if $\delta<\beta\xi$. This step does not impose a length condition, so it is a valid upper bound for the exceptional length mass as well.

For the diameter upper bound put $r=m^{3/4+\xi}$. If $\sum_{i\le k}D_i>r$, one of the two halves has diameter sum exceeding $r/2$. The diameter tail and the same truncation argument as in Lemma~\ref{R:exact:concentration} bound this by $Ckr^{-\beta}$. Pay for the exact length using the independent other half. Summing over $k\le m^{\alpha+\delta}$ gives
\[
 C r^{-\beta}\sum_{k\le m^{\alpha+\delta}}k^{1-1/\alpha}
 \le C m^{\alpha-1-\beta\xi+\delta(2-1/\alpha)}
 =o(V_m)
\]
for sufficiently small $\delta$. Independence was used between disjoint sets of pieces, never between the length and diameter of a piece. Finally $d_0\sum H_i\le A\le D\le\sum D_i+O(1)$. Apply the argument with slightly smaller tolerances to absorb fixed geometric constants.
\end{proof}

\subsection{Critical bridges at a fixed height}
This conditioning requires a different denominator, $B_h$, but only a simpler part of the preceding local argument.
\begin{theorem}\label{R:exact:height}
Under critical weights normalized on strict bridges of exact height $h$ with free terminal port,
\[
 L=h^{4/3+o(1)},\qquad D=h^{1+o(1)}
\]
in probability as every integer $h\to\infty$. The same holds after restricting to any event whose critical mass is at least a fixed positive fraction of $B_h$.
\end{theorem}
\begin{proof}
The mean length is $M_h/B_h\le Ch^{4/3}$, proving the upper length tail at every positive exponent slack. The height variable has Laplace deficit comparable to $s^\beta$ and tail $O(x^{-\beta})$. Its support is an unbounded subset of the integers. The modulus-one characters therefore form a finite set on the Fourier torus; the same near-origin estimate translates to each of them. Lemma~\ref{R:exact:concentration}'s Fourier proof gives
$\sup_j\Prob(\sum_{i\le k}H_i=j)\le Ck^{-1/\beta}$.
Splitting into halves as above shows
\[
 \sum_{k\le h^{\beta-\delta}}\Prob\{\textstyle\sum_{i\le k}H_i=h\}
 \le Ch^{-\beta}\left(1+\sum_{2\le k\le h^{\beta-\delta}}k^{1-1/\beta}\right)
 =o(h^{\beta-1})=o(B_h).
\]
For $k>h^{\beta-\delta}$, length below $h^{4/3-\epsilon}$ costs at most
$\exp\{1-c h^{\alpha\epsilon-\delta}\}$ by the length deficit. There are at most $h$ possible counts, so choosing $\delta<\alpha\epsilon$ proves the lower length tail. The diameter is at least $d_0h$; its upper tail at $h^{1+\epsilon}$ has relative mass at most $Ch^{-\epsilon/4}$ by Theorem~\ref{R:renew:finite-input}. A restriction retaining a fixed fraction of the denominator multiplies all these vanishing upper bounds by only a constant.
\end{proof}
